% Insertar bajo el encabezado 2.7 Seguridad.
% Este archivo incluye los encabezados 2.7.1, 2.7.2 y 2.7.3.

\subsubsection{Autenticación y gestión de sesiones}

El acceso al panel web se realiza mediante correo electrónico y contraseña
a través de \texttt{/api/login}. El servidor consulta la cuenta y exige que
se encuentre activa y no esté eliminada. La contraseña recibida se compara
con el hash almacenado mediante bcrypt; las contraseñas no se almacenan en
texto plano. Cuando la cuenta no cumple las condiciones o la contraseña no
coincide, se devuelve una respuesta HTTP 401 con un mensaje genérico, sin
indicar cuál de las credenciales es incorrecta. La ruta de inicio de sesión
limita las solicitudes a diez intentos por minuto mediante el middleware
de control de frecuencia; al exceder el límite responde con HTTP 429.

Después de una autenticación válida, el servidor emite un token JWT firmado
con el algoritmo HS256 y una vigencia de ocho horas. El token identifica al
usuario e incluye la versión de su sesión. En cada petición protegida, el
middleware comprueba la firma y la expiración, consulta nuevamente la cuenta
y compara la versión del token con \texttt{version\_sesion}. También verifica
que el usuario permanezca activo y pertenezca a la clínica solicitada.

La interfaz conserva el token en \texttt{sessionStorage}, bajo la clave
\texttt{clinica-token}. El interceptor HTTP lo incorpora al encabezado
\texttt{Authorization} utilizando el esquema Bearer. Este almacenamiento
está asociado a la sesión de la pestaña y no se utiliza
\texttt{localStorage} para conservar el token. Sin embargo,
\texttt{sessionStorage} es accesible desde JavaScript: no equivale a una
cookie \texttt{HttpOnly} ni proporciona protección por sí solo frente a
la ejecución de scripts maliciosos.

La implementación no utiliza tokens de refresco ni renovación automática.
Cuando el JWT vence, la API responde con HTTP 401; el interceptor elimina
el token almacenado y dirige al usuario a la pantalla de acceso para que
inicie sesión nuevamente. Esta conducta también se aplica cuando el
servidor detecta una sesión invalidada.

El cierre de sesión llama a \texttt{/api/logout}. El servidor incrementa
\texttt{version\_sesion}, de modo que los tokens previamente emitidos para
esa cuenta dejan de ser válidos en las siguientes solicitudes, incluso
si todavía no han alcanzado su fecha de expiración. Después de recibir la
confirmación, la interfaz limpia el almacenamiento de sesión y vuelve al
inicio de sesión. Esta invalidación afecta a las sesiones anteriores de
la cuenta, no solamente a la pestaña que solicita el cierre. El cambio y
el restablecimiento de contraseña también incrementan la versión de sesión.

La clave de firma del JWT y las credenciales de base de datos y almacenamiento
se obtienen desde variables de entorno del backend. El archivo local
\texttt{backend/.env} está excluido del repositorio y estas credenciales
no se incorporan al código del frontend.

\subsubsection{Autorización por roles o perfiles}

La autorización se aplica en dos niveles. En la interfaz se muestran las
opciones y acciones disponibles según la sesión. En el servidor se
comprueban los permisos de cada operación y el alcance de los registros
solicitados. La comprobación del backend es necesaria aunque se invoque
directamente una ruta HTTP sin utilizar el menú de la aplicación.

Las asociaciones de usuarios y roles se conservan en
\texttt{roles\_modelos}; los permisos de cada rol se almacenan en
\texttt{permisos\_roles}, y los permisos directos de usuarios en
\texttt{permisos\_modelos}. El servidor recupera esta información al
validar las solicitudes. Los permisos distinguen operaciones como lectura,
creación, actualización y eliminación. Los perfiles internos tienen las
siguientes responsabilidades:

\begin{itemize}
  \item \textbf{Administrador:} administra cuentas, roles y permisos y
  supervisa los registros del sistema dentro de las clínicas a las que
  pertenece. El rol interno \texttt{Super Admin} satisface las comprobaciones
  generales de permisos; las rutas de administración utilizan además la
  comprobación \texttt{requireAdmin}.

  \item \textbf{Odontólogo:} gestiona los pacientes que tiene asignados
  y sus solicitudes radiográficas, de acuerdo con los permisos otorgados.
  Las consultas de pacientes y radiografías comprueban la relación de
  asignación, y las consultas de citas y recetas incorporan también el
  identificador del profesional. No dispone de administración de cuentas
  y permisos por el hecho de tener el rol de odontólogo.

  \item \textbf{Recepcionista:} consulta, registra y actualiza datos
  administrativos de pacientes; consulta odontólogos y horarios; gestiona
  citas y asigna pacientes a profesionales. El servidor limita los campos
  visibles y modificables y no le concede acceso a radiografías,
  antecedentes clínicos, recetas o administración de cuentas. La
  revocación de una asignación está reservada al administrador.
\end{itemize}

CERPAX no es un perfil con inicio de sesión en el panel. Su acceso se limita
al portal de entrega mediante un token temporal vinculado a una solicitud.
La validación exige que el enlace esté vigente, no haya sido revocado ni
utilizado y corresponda a una solicitud pendiente. El token se conserva
como hash en la base de datos. Un enlace inválido o no disponible genera
HTTP 410. El paciente se representa como registro clínico y no como actor
de acceso al panel descrito en esta solución.

Para el aislamiento entre clínicas, el servidor verifica que el valor del
encabezado \texttt{X-Company-ID} pertenezca a las clínicas autorizadas del
usuario. Posteriormente incorpora ese alcance en las consultas. Los
permisos insuficientes generan HTTP 403; algunos accesos a registros fuera
del alcance autorizado responden con HTTP 404 para no revelar su existencia.
El acceso local de desarrollo está condicionado a la configuración y a
la dirección de conexión; se deshabilita cuando el entorno es producción.

\subsubsection{Validación de entradas y protección}

La interfaz utiliza controles de formulario para orientar la captura de
datos, como campos obligatorios, tipos de entrada y opciones disponibles.
El backend realiza su propia validación, ya que los datos pueden enviarse
sin pasar por la interfaz. Las reglas comprueban campos permitidos,
identificadores numéricos, formatos de correo, longitudes máximas,
valores de listas, números finitos y estructuras JSON. Las validaciones
clínicas comprueban fechas y horarios y rechazan intervalos inconsistentes
y superposiciones de citas.

Las operaciones con PostgreSQL utilizan consultas parametrizadas: los
valores recibidos se envían separados del texto SQL mediante marcadores
como \texttt{\$1} y \texttt{\$2}. Los nombres de módulos y campos se
resuelven desde el catálogo y los metadatos permitidos, y los identificadores
SQL se validan antes de incorporarse a las consultas dinámicas. Estas
medidas reducen el riesgo de inyección SQL sin atribuir a una consulta
parametrizada la validación de nombres de tablas o columnas.

Para las radiografías se admiten PDF, PNG, JPEG y WebP con un tamaño máximo
de \mbox{50 MiB}. En la carga directa a Supabase Storage, la API valida el
nombre, la extensión, el tipo MIME y el tamaño declarado antes de emitir
la URL firmada. Al confirmar, comprueba la existencia del objeto, su tamaño
y la correspondencia de su ruta con el acceso temporal. En las rutas de
carga multipart que reciben el archivo en el servidor se utiliza Multer
para limitar el tamaño y se identifica el formato a partir de la firma
binaria. La validación de metadatos de la carga directa no se presenta como
una inspección completa del contenido del archivo.

Los objetos radiográficos se mantienen en almacenamiento privado. Antes
de abrir un archivo, la API vuelve a comprobar la autorización del usuario
sobre la solicitud. Las operaciones que completan una entrega utilizan
transacciones y restricciones de unicidad para conservar la coherencia
entre la solicitud, el archivo y el acceso utilizado.

Express aplica los encabezados de seguridad de Helmet y una política de
seguridad de contenido que restringe los orígenes de scripts, imágenes y
conexiones. La API comprueba el encabezado de origen cuando está presente
y rechaza los orígenes no permitidos con HTTP 403. Este control complementa
la autenticación y no la sustituye. El procesamiento de cuerpos JSON está
limitado a \mbox{2 MiB}, de forma independiente del límite de archivos.

El manejador centralizado de errores devuelve respuestas JSON y evita
exponer consultas internas, contraseñas o cadenas de conexión. Se utiliza
HTTP 400 para errores de formato de la petición, HTTP 401 para autenticación,
HTTP 403 para autorización, HTTP 409 para conflictos y HTTP 422 para
validaciones de datos, según la ruta y la condición detectada. Los errores
internos generan un mensaje genérico con HTTP 500. Las respuestas de la
API incluyen \texttt{Cache-Control: no-store} para evitar su almacenamiento
en cachés HTTP.

Estas medidas se respaldan con pruebas del backend que comprueban ausencia
de sesión, permisos insuficientes, acceso entre clínicas, pacientes no
asignados, enlaces vencidos o revocados y validación de archivos y datos.
Los resultados de cada escenario deben acompañarse de las evidencias
correspondientes en la sección de pruebas.
